% from ega1-9-fr.tex
\oldpage[I]{170}
En effet, en vertu de (\hyperref[I.1.6.5-fr]{1.6.5}),
$\mathscr{F}\otimes_S\mathscr{G}$ est canoniquement isomorphe au faisceau
associé au $(B\otimes_A C)$-module
\[
  \bigl(M\otimes_B(B\otimes_A C)\bigr)
  \otimes_{B\otimes_A C}
  \bigl((B\otimes_A C)\otimes_C N\bigr)
\]
et en raison des isomorphismes canoniques entre produits tensoriels, ce
dernier est isomorphe à
\[
  M\otimes_B(B\otimes_A C)\otimes_C N
  =(M\otimes_B B)\otimes_A(C\otimes_C N)=M\otimes_A N.
\]
\end{proof}

\begin{proposition}[9.1.4]
\label{I.9.1.4-fr}
Soient $f:T\to X$, $g:T\to Y$ deux $S$-morphismes, $\mathscr{F}$ (resp.
$\mathscr{G}$) un $\mathscr{O}_X$-Module (resp. $\mathscr{O}_Y$-Module)
quasi-cohérent. On a alors
\[
  (f,g)_S^*(\mathscr{F}\otimes_S\mathscr{G})
  =f^*(\mathscr{F})\otimes_{\mathscr{O}_T}g^*(\mathscr{G}).
\]
\end{proposition}

\begin{proof}
Si $p$, $q$ sont les projections de $X\times_S Y$, la formule résulte en
effet des relations $(f,g)_S^*\circ p^*=f^*$ et
$(f,g)_S^*\circ q^*=g^*$ (\hyperref[0.3.5.5-fr]{0, 3.5.5}), et du fait
qu'une image réciproque d'un produit tensoriel de faisceaux algébriques est
le produit tensoriel de leurs images réciproques
(\hyperref[0.4.3.3-fr]{0, 4.3.3}).
\end{proof}

\begin{corollary}[9.1.5]
\label{I.9.1.5-fr}
Soient $f:X\to X'$, $g:Y\to Y'$ deux $S$-morphismes, $\mathscr{F}'$ (resp.
$\mathscr{G}'$) un $\mathscr{O}_{X'}$-Module (resp.
$\mathscr{O}_{Y'}$-Module) quasi-cohérent. On a alors
\[
  (f\times_S g)^*(\mathscr{F}'\otimes_S\mathscr{G}')
  =f^*(\mathscr{F}')\otimes_S g^*(\mathscr{G}').
\]
\end{corollary}

\begin{proof}
Cela résulte de (\hyperref[I.9.1.4-fr]{9.1.4}) et du fait que
$f\times_S g=(f\circ p,g\circ q)_S$, $p$ et $q$ étant les projections de
$X\times_S Y$.
\end{proof}

\begin{corollary}[9.1.6]
\label{I.9.1.6-fr}
Soient $X$, $Y$, $Z$ trois $S$-préschémas, $\mathscr{F}$ (resp.
$\mathscr{G}$, $\mathscr{H}$) un $\mathscr{O}_X$-Module (resp.
$\mathscr{O}_Y$-Module, $\mathscr{O}_Z$-Module) quasi-cohérent~; le faisceau
$\mathscr{F}\otimes_S\mathscr{G}\otimes_S\mathscr{H}$ est l'image
réciproque de
$(\mathscr{F}\otimes_S\mathscr{G})\otimes_S\mathscr{H}$ par l'isomorphisme
canonique de $X\times_S Y\times_S Z$ sur
$(X\times_S Y)\times_S Z$.
\end{corollary}

\begin{proof}
En effet, cet isomorphisme s'écrit $(p_1,p_2)_S\times_S p_3$, en désignant
par $p_1$, $p_2$, $p_3$ les projections de $X\times_S Y\times_S Z$.

De même, l'image réciproque de $\mathscr{G}\otimes_S\mathscr{F}$ par
l'isomorphisme canonique de $X\times_S Y$ sur $Y\times_S X$ est
$\mathscr{F}\otimes_S\mathscr{G}$.
\end{proof}

\begin{corollary}[9.1.7]
\label{I.9.1.7-fr}
Si $X$ est un $S$-préschéma, tout $\mathscr{O}_X$-Module quasi-cohérent
$\mathscr{F}$ est l'image réciproque de
$\mathscr{F}\otimes_S\mathscr{O}_S$ par l'isomorphisme canonique de $X$ sur
$X\times_S S$ (\hyperref[I.3.3.3-fr]{3.3.3}).
\end{corollary}

\begin{proof}
En effet, cet isomorphisme est $(1_X,\varphi)_S$, où $\varphi$ est le
morphisme structural $X\to S$, et le corollaire résulte de
(\hyperref[I.9.1.4-fr]{9.1.4}) et du fait que
$\varphi^*(\mathscr{O}_S)=\mathscr{O}_X$.
\end{proof}

\begin{env}[9.1.8]
\label{I.9.1.8-fr}
Soient $X$ un $S$-préschéma, $\mathscr{F}$ un $\mathscr{O}_X$-Module
quasi-cohérent, $\varphi:S'\to S$ un morphisme~; on désigne par
$\mathscr{F}_{(\varphi)}$ ou $\mathscr{F}_{(S')}$ le faisceau
quasi-cohérent $\mathscr{F}\otimes_S\mathscr{O}_{S'}$, sur
$X\times_S S'=X_{(\varphi)}=X_{(S')}$~; donc
$\mathscr{F}_{(S')}=p^*(\mathscr{F})$, où $p$ est la projection
$X_{(S')}\to X$.
\end{env}

\begin{proposition}[9.1.9]
\label{I.9.1.9-fr}
Soit $\varphi':S''\to S'$ un morphisme. Pour tout
$\mathscr{O}_X$-Module quasi-cohérent $\mathscr{F}$ sur le $S$-préschéma
$X$, $(\mathscr{F}_{(\varphi)})_{(\varphi')}$ est l'image réciproque de
$\mathscr{F}_{(\varphi\circ\varphi')}$ par l'isomorphisme canonique
$(X_{(\varphi)})_{(\varphi')}\xrightarrow{\sim}
X_{(\varphi\circ\varphi')}$ (\hyperref[I.3.3.9-fr]{3.3.9}).
\end{proposition}

\begin{proof}
Cela résulte aussitôt des définitions et de
(\hyperref[I.3.3.9-fr]{3.3.9}), et s'écrit encore
\[
  (\mathscr{F}\otimes_S\mathscr{O}_{S'})
  \otimes_{S'}\mathscr{O}_{S''}
  =\mathscr{F}\otimes_S\mathscr{O}_{S''}.
  \tag{9.1.9.1}\label{I.9.1.9.1-fr}
\]
\end{proof}

\begin{proposition}[9.1.10]
\label{I.9.1.10-fr}
Soient $Y$ un $S$-préschéma, $f:X\to Y$ un $S$-morphisme. Pour tout
$\mathscr{O}_Y$-Module quasi-cohérent $\mathscr{G}$ et tout morphisme
$S'\to S$, on a
\[
  (f_{(S')})^*(\mathscr{G}_{(S')})
  =(f^*(\mathscr{G}))_{(S')}.
\]
\end{proposition}

\begin{proof}
Cela résulte aussitôt de la commutativité du diagramme
